\section{Algoritma Euclid dan Euclid Diperluas}

\subsection{Algoritma Euclid}
Algoritma Euclid menghitung $\gcd(a, b)$ secara efisien. Ide: $\gcd(a, b) = \gcd(b, a \bmod b)$. Proses berulang hingga $b = 0$, maka $\gcd = a$.

\begin{table}[!htbp]
\centering
\caption{Contoh: $\gcd(48, 18)$}
\label{tab:euclid}
\begin{tabular}{@{}ccc@{}}
\toprule
$a$ & $b$ & $a \bmod b$ \\
\midrule
48 & 18 & 12 \\
18 & 12 & 6 \\
12 & 6 & 0 \\
6 & 0 & --- \\
\bottomrule
\end{tabular}
\end{table}
Jadi $\gcd(48, 18) = 6$.

\subsection{Algoritma Euclid Diperluas}
Algoritma Euclid diperluas tidak hanya menghitung $\gcd(a, b)$ tetapi juga koefisien $x$ dan $y$ sehingga $ax + by = \gcd(a, b)$. Koefisien $x$ memberikan invers modular dari $a$ modulo $b$ jika $\gcd(a, b) = 1$.

\begin{lstlisting}[style=matlab,caption={Algoritma Euclid Diperluas dalam MATLAB}]
function [g, x, y] = extended_gcd(a, b)
    if b == 0
        g = a; x = 1; y = 0;
    else
        [g, x1, y1] = extended_gcd(b, mod(a, b));
        x = y1;
        y = x1 - floor(a/b) * y1;
    end
end
% Contoh: invers 7 mod 26 (untuk Caesar)
% extended_gcd(7, 26) -> x=15, jadi 7^(-1) mod 26 = 15
\end{lstlisting}

Implementasi lengkap dalam C, C++, dan Pascal tersedia di Bab 14.
